
% =====================================================================
\section{Cosmology II: can the superfluid be the dark matter?}
\label{sec:dm}
% =====================================================================

The second part of hypothesis H3 concerns the dark matter. Here the fermions are taken to be heavy
enough to be non-relativistic matter today, and the question is whether the paired state of
Secs.~\ref{sec:cjt} and \ref{sec:order} can account for what is observed. The observations used are
of three kinds: the amount of small-scale structure in the universe, measured by the absorption
lines of intergalactic hydrogen, the Lyman-$\alpha$ forest; the densities and velocities of the dark
matter in galaxies and clusters; and the limit on the scattering of dark matter by dark matter from
colliding clusters of galaxies.

\subsection{The superfluid as a cosmological fluid}
\label{sec:sffluid}

Suppose that the dark matter consists of the fermions of the model: an excess of particles over
antiparticles, whose size is a free number that is chosen to give the observed density. At zero
temperature they are in the paired state. Its response to slow perturbations of long wavelength is
that of a fluid with a pressure $P(\rho)$ that depends on the mass density $\rho$ only: perturbations
of the density propagate as first sound, with the speed $c_s$ given by $c_s^2=dP/d\rho$
(Sec.~\ref{sec:sound}). Two limits bracket the possibilities.

\emph{Weak pairing}, pairs that overlap. The pressure is the degeneracy pressure of the fermions,
which pairing lowers only by terms of the relative order $(\Delta/\EF)^2$. With $\nu=8$ internal
states the Fermi momentum at the number density $n$ is $\kF=(6\pi^2n/\nu)^{1/3}$, and with
$x=\kF/m$ the sound speed is
\begin{equation}
c_s^2=\frac{x^2}{3\,(1+x^2)},
\label{eq:csdeg}
\end{equation}
which is $v_F^2/3$ for non-relativistic fermions with the Fermi velocity $v_F=\kF/m$. At the mean
cosmic density of dark matter the Fermi velocity today is fixed by the mass alone. It grows as $1/a$
towards the past, where $a$ is the scale factor of the universe, and the gas was relativistic at
early times, which the program includes exactly.

\emph{Bound pairs}, or larger bound clusters, that have condensed. If their mutual interaction is
repulsive, the condensate has the equation of state of a self-interacting scalar field
\cite{ColpiShapiroWasserman1986,LiRindlerDallerShapiro2014},
\begin{equation}
P=\frac{\rho_0}{3}\Big[\sqrt{1+\rho/\rho_0}-1\Big]^2,
\label{eq:csw}
\end{equation}
with one constant $\rho_0$ that measures the strength of the repulsion. At low density this is
$P=\rho^2/12\rho_0$, the equation of state of a condensate with the sound speed of Bogoliubov
\cite{Bogoliubov1947}, $c_s^2=\rho/6\rho_0$; at high density it is that of radiation, $P=\rho/3$.
The condensate is characterized by one number, the sound speed $c_{s0}$ today at the mean density.
A self-gravitating core that is supported by this pressure has the \emph{Thomas-Fermi radius}
\begin{equation}
R_{TF}=\frac{\pi\,c_s}{\sqrt{4\pi G\rho}},
\label{eq:RTF}
\end{equation}
which does not depend on the density, because $c_s^2\propto\rho$; $G$ is the constant of
gravitation. Cores of this kind, with radii of the order of a kiloparsec, are the motivation of the
proposals of superfluid dark matter \cite{BerezhianiKhoury2015}.

In both limits the pressure is the \emph{smallest} that the constituents can have; a thermal
population of the same particles has more. The bounds derived below are therefore lower limits on
what any state of these particles must satisfy.

\subsection{Simulation}

The cold dark matter of the Einstein-Boltzmann solver of Sec.~\ref{sec:solver} was replaced by a
fluid with the equation of state and the sound speed of either limit. Tests of the program verify
the conservation of energy, the relation $c_s^2=dP/d\rho$ and the limits of both equations of state,
and that a fluid with a negligible sound speed reproduces cold dark matter to $2\times10^{-4}$. The
linear power spectrum of the matter today was computed for wavenumbers up to 250 Mpc$^{-1}$ and
divided by that of cold dark matter (Fig.~\ref{fig:sfdm}). The \emph{half mode} is the wavenumber at
which this ratio first falls to $1/4$. It is compared with the half mode of \emph{thermal warm dark
matter}, particles that were once in thermal equilibrium and whose residual velocities erase
small-scale structure, at the lower limits on their mass from the Lyman-$\alpha$ forest: 5.3 keV
\cite{Irsic2017} and, more conservatively, 3.5 keV; the relation between that mass and the half mode
is the fitting formula of Ref.~\cite{Viel2005}. Table~\ref{tab:sfdm} lists the results and
Table~\ref{tab:sfdmbounds} the bounds that follow. A more recent analysis of the Lyman-$\alpha$
forest gives 5.7 keV \cite{Irsic2024}, which would strengthen the bounds slightly.

\begin{figure*}
\includegraphics[width=\textwidth]{fig_sfdm_transfer}
\caption{\label{fig:sfdm}Linear power spectrum of the matter, relative to cold dark matter, if the
dark matter is the superfluid. (a) Weak pairing, for several fermion masses. (b) Condensate of bound
bosons, for several sound speeds $c_{s0}$ today at the mean density; the largest corresponds to a
Thomas-Fermi radius of 1 kpc. The vertical lines mark the half modes of thermal relics of 3.5 and
5.3 keV. The oscillations are sound waves of the dark fluid.}
\end{figure*}

\begin{table*}
\caption{\label{tab:sfdm}Superfluid dark matter in the two limits of Sec.~\ref{sec:sffluid}. $c_{s0}$:
sound speed today at the mean cosmic density, in units of the speed of light. Half mode: wavenumber,
in $h$ Mpc$^{-1}$, at which the power falls to $1/4$ of that of cold dark matter. Relic: mass of the
thermal warm-dark-matter particle with the same half mode. $R_{TF}$: Thomas-Fermi radius,
Eq.~(\ref{eq:RTF}).}
\begin{ruledtabular}
\begin{tabular}{lcccc}
model & $c_{s0}/c$ & half mode ($h$ Mpc$^{-1}$) & relic (keV) & $R_{TF}$ \\
\hline
{{ROWS:sfdm_table|cols=0,1,3-5}}
\end{tabular}
\end{ruledtabular}
\end{table*}

\begin{table*}
\caption{\label{tab:sfdmbounds}Bounds on superfluid dark matter from the requirement that its half mode
is not smaller than that of thermal warm dark matter at the limits from the Lyman-$\alpha$ forest.
$v_{F0}$: Fermi velocity today at the mean density.}
\begin{ruledtabular}
\begin{tabular}{ccccccc}
relic (keV) & its half mode ($h$ Mpc$^{-1}$) & fermion mass $>$ (keV) & $v_{F0}<$ (m/s) &
condensate: $c_{s0}/c<$ & $c_{s0}<$ (mm/s) & $R_{TF}<$ (pc) \\
\hline
{{ROWS:sfdm_bounds_table}}
\end{tabular}
\end{ruledtabular}
\end{table*}

\subsection{What the bounds imply}

\emph{Weak pairing.} The fermion mass must exceed {{GEN:sfdm_mmin_5p3}} keV, or
{{GEN:sfdm_mmin_3p5}} keV for the conservative benchmark. A paired state requires a degenerate gas,
because its critical temperature is a fraction of the Fermi energy (Sec.~\ref{sec:thermo}). A gas of
the mass density $\rho$ whose particles move with the velocity dispersion $\sigma$ is degenerate if
its Fermi velocity is at least $\sigma$, which is the case only for fermions lighter than
\begin{equation}
m_{\rm deg}=\Big[\frac{6\pi^2\hbar^3\rho}{\nu\,\sigma^3}\Big]^{1/4}.
\label{eq:mdeg}
\end{equation}
This is the bound of Tremaine and Gunn \cite{TremaineGunn1979} for this model. Table~\ref{tab:halos}
evaluates it for three observed systems.

\begin{table*}
\caption{\label{tab:halos}Could the particles of the model be degenerate, or their bound pairs
condensed, in observed systems of dark matter? $\rho$: density; $\sigma$: velocity dispersion;
$m_{\rm deg}$: Eq.~(\ref{eq:mdeg}); $m_{B,\max}$: Eq.~(\ref{eq:mbmax}). The last two columns give,
for fermions of 1.5 keV, the Fermi velocity at the density $\rho$ and the factor by which the density
would have to be larger for them to be degenerate.}
\begin{ruledtabular}
\begin{tabular}{lccccccc}
system & $\rho$ ($M_\odot$ pc$^{-3}$) & $\rho$ (GeV cm$^{-3}$) & $\sigma$ (km/s) & $m_{\rm deg}$ (keV) &
$m_{B,\max}$ (eV) & $v_F$ (km/s) & $\rho_{\rm deg}/\rho$ \\
\hline
{{ROWS:dm_halo_table}}
\end{tabular}
\end{ruledtabular}
\end{table*}

In the densest systems of dark matter that are observed, the cores of dwarf spheroidal galaxies,
fermions with eight internal states are degenerate only below {{GEN:dm_mdeg_dwarf}} keV, whereas
the formation of structure demands more than {{GEN:sfdm_mmin_5p3}} keV. \emph{In the regime of weak
pairing the dark matter in galaxies is a classical gas, not a superfluid.}

\emph{Bound pairs.} The sound speed of the condensate today must be below
{{GEN:sfdm_csmax_5p3}} of the speed of light, and the Thomas-Fermi radius below
{{GEN:sfdm_rtf_pc_5p3}} pc. The reason for so strong a bound is visible in Fig.~\ref{fig:sfdm}(b):
such a condensate is like radiation in the early universe, and all perturbations that enter the
horizon during that phase are erased. \emph{The pressure of the condensate cannot produce the cores of
kiloparsec size that motivate superfluid dark matter;} for a core of 1 kpc the half mode would be at
{{GEN:sfdm_khalf_1kpc}} $h$ Mpc$^{-1}$, more than a hundred times too small. There is a second,
independent obstacle. Bosons of the mass $m_B$ are condensed at the density $\rho$ and the
temperature that corresponds to the velocity dispersion $\sigma$, $k_BT=m_B\sigma^2$, only if the
number of particles in a cube with the edge of the thermal wavelength $\lambda$ satisfies
$n\lambda^3\ge2.612$, that is, only if they are lighter than
\begin{equation}
m_{B,\max}=\Big[\frac{(2\pi)^{3/2}\hbar^3\rho}{2.612\,\sigma^3}\Big]^{1/4},
\label{eq:mbmax}
\end{equation}
which is {{GEN:dm_mbmax_dwarf}} eV in a dwarf galaxy (Table~\ref{tab:halos}). The limit on their
self-interaction demands the opposite, as shown next.

\emph{Self-interaction.} Colliding clusters of galaxies limit the cross section per unit mass of
dark matter to about 1 cm$^2$/g \cite{Markevitch2004,Randall2008}. A pair that is bound by the
Yukawa potential of the radial mode has a size, and a scattering length, of at least the range
$1/m_s$ of that potential, and binding requires $\DVS>D_c$ (Sec.~\ref{sec:dimer}). Identical
bosons with the scattering length $a_s$ scatter at low energy with the cross section $8\pi a_s^2$;
with $a_s=1/m_s$ and the mass $2m$ of a pair, the two conditions together give a smallest fermion
mass, $m\ge[4\pi D_c^2/(\aVS^2\kappa)]^{1/3}$ with $\kappa=1$ cm$^2$/g, listed in
Table~\ref{tab:selfint} for several couplings.

\begin{table}
\caption{\label{tab:selfint}Smallest fermion mass for which a gas of bound pairs satisfies the limit of
1 cm$^2$/g on the self-interaction of dark matter, the corresponding mass and range of the radial
mode, and the phase-space density $n\lambda^3$ of the pairs in the core of a dwarf spheroidal galaxy;
condensation would need $n\lambda^3\ge2.612$.}
\begin{ruledtabular}
\begin{tabular}{ccccc}
$\aVS$ & $m$ (MeV) & $m_s$ (MeV) & $1/m_s$ (fm) & $n\lambda^3$ \\
\hline
{{ROWS:dm_selfint_table}}
\end{tabular}
\end{ruledtabular}
\end{table}

Bound pairs are acceptable dark matter only if the fermions are heavier than about
{{GEN:dm_mmin_alpha1}} MeV, for $\aVS=1$, with a radial mode of nuclear range. At such masses the
phase-space density in halos is smaller than the value needed for condensation by more than twenty
orders of magnitude.

\emph{Bound clusters.} With eight internal states and an attraction between all of them, up to eight
fermions can occupy the same orbital, and bound states of more than two fermions are expected, as
nuclei are bound states of nucleons. Such ``nuggets'' of fermions bound by a scalar field have been
studied as dark matter \cite{WiseZhang2014,GreshamLouZurek2017}. Their interior is degenerate and,
by Sec.~\ref{sec:cjt}, paired: the superfluid would exist \emph{inside} the particles of dark matter,
as it does inside nuclei and neutron stars, without any effect on galactic scales.
Section~\ref{sec:normal} gave the limit of large clusters: matter is bound by itself if
$C=\gVS^2m^2/m_s^2$ exceeds {{GEN:cjt_c_crit}}, or {{GEN:cjt_c_crit_sea}} if the energy of the
negative-energy states is included, and it is then in a state in which the Vorton-Scalar has no
expectation value and the fermions are massless. For the example of Table~\ref{tab:selfint} with
$\aVS=1$ the coupling is $C=35.5$, and a fermion inside bulk matter weighs {{GEN:cjt_ebind_355}} per
cent less than a free one, or {{GEN:cjt_ebind_355_sea}} per cent with the negative-energy states.
Whether such clusters form, and with what abundance, is outside this work.

\subsection{Verdict on dark matter}

\begin{figure*}
\includegraphics[width=\textwidth]{fig_dm_windows}
\caption{\label{fig:dmwindows}The fermion masses that a superfluid on galactic scales needs and those
that the observations allow. For weak pairing the fermions must be degenerate in halos, which requires
a mass below the value of Eq.~(\ref{eq:mdeg}), while the Lyman-$\alpha$ forest requires a larger one.
For bound pairs, condensation in a halo requires light pairs, while the limit on the self-interaction
requires heavy ones (ticks: other values of the coupling). In neither case do the ranges overlap.}
\end{figure*}

The hypothesis has two readings, and they have different answers. Figure~\ref{fig:dmwindows}
summarizes the numbers on which the first answer rests.

\emph{The superfluid state explains the observed properties of dark matter,} as in the proposals in
which a condensate of galactic size shapes rotation curves or cores. \emph{No.} The simulations
exclude both ways in which the condensate could act on galactic scales: degenerate fermions that are
light enough to be superfluid in halos erase the small-scale structure, and a condensate of bound
pairs with enough pressure to support cores does so even more strongly.

\emph{The dark matter is made of the particles whose condensate is the superfluid.} \emph{Possible,
and untested by these observations.} Heavy bound pairs or larger clusters, with an abundance that is
set by an unexplained excess of particles over antiparticles, behave as cold dark matter with a
self-interaction of nuclear size. Nothing in the present data distinguishes them from other cold dark
matter, and the model does not predict their abundance.

\subsection{What has been established, and what has not}

Established: the suppression of structure by superfluid dark matter in both limits and the resulting
bounds; and the elementary conditions for degeneracy, for condensation and for an acceptable
self-interaction in halos, which the masses allowed by those bounds violate by large factors.

Not established, or assumed: the translation of the limits from the Lyman-$\alpha$ forest into a
half mode uses a fitting formula for thermal relics, whereas the spectrum of a fluid with sound
waves has a different shape beyond the half mode; a dedicated analysis of those data could shift the
bounds by a factor of order one, not by the orders of magnitude that separate them from the
interesting values. For bound pairs a repulsive interaction between the pairs was assumed; with
eight internal states it may be attractive, in which case the condensate is unstable to the
formation of larger clusters. The sound speed of the degenerate gas was that of free fermions; the
correction from the mean field of the Vorton-Scalar is of the relative order $\kF|a_s|$, with $a_s$
the scattering length, which the limit on the self-interaction keeps below {{GEN:cjt_kfa_dwarf}} in
a dwarf galaxy for fermions of 1 keV. The abundance of the dark matter, the origin of the excess of
particles, the thermal history of the dark sector and the formation of bound states were not
addressed.
